% ============================================================
%  Supplementary Methods: Search Strategies, Eligibility Criteria,
%  and Risk of Bias Assessment
% ============================================================

\section*{Supplementary Methods}
\label{app:methods}

This supplementary file provides complete methodological details referenced in Section~\ref{sec:methods} of the main text, including database-specific search strategies, detailed eligibility criteria with exclusion codes, and the adapted QUADAS-2 framework for risk of bias assessment.

% ----------------------------------------------------------------
% Search Strategies
% ----------------------------------------------------------------
\section{Database-Specific Search Strategies}
\label{app:methods:search}

Searches were executed in three rounds: January~2024 (covering 2017--2023), June~2024 (covering 2017--mid-2024), and March~2026 (final update). All searches combined three concept blocks with Boolean AND: (1) histopathology terms, (2) fairness and bias terms, and (3) AI or deep learning terms.

\subsection*{PubMed}

\begin{verbatim}
("histopathology"[MeSH Terms] OR "histopathology"[All Fields] OR 
 "whole slide image"[All Fields] OR "WSI"[All Fields] OR 
 "digital pathology"[All Fields] OR "computational pathology"[All Fields] OR
 "tissue microarray"[All Fields] OR "TMA"[All Fields] OR 
 "FFPE"[All Fields] OR "frozen section"[All Fields] OR 
 "immunohistochemistry"[All Fields] OR "IHC"[All Fields])
AND
("fairness"[All Fields] OR "bias"[MeSH Terms] OR "bias"[All Fields] OR 
 "health equity"[All Fields] OR "algorithmic bias"[All Fields] OR 
 "domain shift"[All Fields] OR "batch effect"[All Fields] OR 
 "selection bias"[All Fields] OR "confounding"[All Fields])
AND
("deep learning"[MeSH Terms] OR "deep learning"[All Fields] OR 
 "machine learning"[MeSH Terms] OR "machine learning"[All Fields] OR 
 "transformer"[All Fields] OR "foundation model"[All Fields] OR 
 "neural network"[All Fields] OR "convolutional neural network"[All Fields] OR
 "self-supervised learning"[All Fields])
\end{verbatim}

Filters: English language, Publication date from 2017/01/01 to 2026/03/31.

\subsection*{Scopus}

\begin{verbatim}
TITLE-ABS-KEY (
  "histopathology" OR "whole slide image" OR "WSI" OR 
  "digital pathology" OR "computational pathology" OR 
  "tissue microarray" OR "TMA" OR "FFPE" OR 
  "frozen section" OR "immunohistochemistry" OR "IHC"
)
AND
TITLE-ABS-KEY (
  "fairness" OR "bias" OR "health equity" OR 
  "algorithmic bias" OR "domain shift" OR "batch effect" OR 
  "selection bias" OR "confounding"
)
AND
TITLE-ABS-KEY (
  "deep learning" OR "machine learning" OR "transformer" OR 
  "foundation model" OR "neural network" OR 
  "convolutional neural network" OR "self-supervised learning"
)
AND (LIMIT-TO (LANGUAGE, "English"))
AND (LIMIT-TO (PUBYEAR, 2017-2026))
\end{verbatim}

\subsection*{Web of Science}

\begin{verbatim}
TS=("histopathology" OR "whole slide image" OR "WSI" OR 
    "digital pathology" OR "computational pathology" OR 
    "tissue microarray" OR "FFPE" OR "frozen section" OR 
    "immunohistochemistry")
AND
TS=("fairness" OR "bias" OR "health equity" OR 
    "algorithmic bias" OR "domain shift" OR "batch effect" OR 
    "selection bias" OR "confounding")
AND
TS=("deep learning" OR "machine learning" OR "transformer" OR 
    "foundation model" OR "neural network" OR 
    "convolutional neural network" OR "self-supervised")
Refined by: English (LANGUAGE) AND 2017-2026 (PUBLICATION YEARS)
\end{verbatim}

\subsection*{IEEE Xplore}

\begin{verbatim}
("Document Content":"histopathology" OR "whole slide image" OR "WSI" OR 
 "digital pathology" OR "computational pathology")
AND
("Document Content":"fairness" OR "bias" OR "health equity" OR 
 "algorithmic bias" OR "domain shift" OR "batch effect")
AND
("Document Content":"deep learning" OR "machine learning" OR 
 "transformer" OR "foundation model" OR "neural network")
Filter: English, Publication Year: 2017-2026
\end{verbatim}

\subsection*{ACM Digital Library}

\begin{verbatim}
{\small
[Title: histopathology OR "whole slide image" OR WSI OR "digital pathology"] 
[Abstract: fairness OR bias OR "health equity" OR "algorithmic bias"] 
[Title: "deep learning" OR "machine learning" OR transformer OR "foundation model"]
}
Filters: English language, 2017-2026
\end{verbatim}

\subsection*{Google Scholar}

Google Scholar searches used a simplified query due to character limits:

\begin{verbatim}
histopathology OR "whole slide image" OR "digital pathology" 
fairness OR bias OR "domain shift" OR "batch effect"
"deep learning" OR "machine learning" OR "foundation model"
\end{verbatim}

Results were screened to the first 200 results per search round (Google Scholar returns up to 1000 results but relevance decreases substantially beyond the first 200).

% ----------------------------------------------------------------
% Eligibility Criteria with Exclusion Codes
% ----------------------------------------------------------------
\section{Eligibility Criteria and Exclusion Codes}
\label{app:methods:eligibility}

\subsection*{Title and Abstract Screening Exclusion Codes (E1--E7)}

During title/abstract screening, records were excluded using the following codes:

\begin{itemize}
  \item[\textbf{E1}] \textbf{Wrong modality}: Study uses non-histopathology imaging (e.g., radiology, dermatoscopy, retinal imaging, cytology smears) without separate histopathology analysis.
  
  \item[\textbf{E2}] \textbf{No fairness/bias analysis}: Study reports model performance but includes no fairness evaluation, subgroup analysis, domain shift assessment, or bias discussion.
  
  \item[\textbf{E3}] \textbf{Review only}: Editorial, commentary, narrative review, systematic review, or position paper without original experimental results.
  
  \item[\textbf{E4}] \textbf{No results}: Conference abstract, workshop proposal, or protocol paper without full experimental results.
  
  \item[\textbf{E5}] \textbf{Non-human subjects}: Study uses only animal models, cell lines without patient association, or synthetic data without clinical validation.
  
  \item[\textbf{E6}] \textbf{Not AI/ML}: Study uses traditional image processing, handcrafted features without machine learning, or purely statistical analysis without predictive modeling.
  
  \item[\textbf{E7}] \textbf{Language}: Full text not available in English.
\end{itemize}

\subsection*{Full-Text Exclusion Codes (F1--F4)}

During full-text assessment, articles were excluded using the following codes:

\begin{itemize}
  \item[\textbf{F1}] \textbf{Non-histopathology} ($n=193$): Upon full-text review, the study was confirmed to use non-histopathology modalities only. This includes radiology (CT, MRI, X-ray), dermatoscopy, fundus imaging, OCT, or cytology without a separate histopathology cohort.
  
  \item[\textbf{F2}] \textbf{No fairness/subgroup analysis} ($n=289$): The study reports overall performance metrics only, with no subgroup analysis, domain shift evaluation, cross-site validation, or bias discussion. Papers mentioning ``bias'' only in the context of statistical bias (bias--variance tradeoff) or model bias terms were excluded.
  
  \item[\textbf{F3}] \textbf{Duplicate cohort} ($n=53$): Multiple papers report on the exact same dataset, model, and experimental results. In such cases, only the most complete or earliest publication was retained to avoid double-counting evidence.
  
  \item[\textbf{F4}] \textbf{Insufficient detail/full text} ($n=22$): Conference abstracts without full text, papers where key methodological details were unavailable, or studies where fairness-related results were mentioned but not quantified.
\end{itemize}

\subsection*{Worked Examples}

\begin{itemize}
  \item \textbf{Example E1}: A study evaluating fairness of a chest X-ray pneumonia detector would be excluded under E1 (wrong modality), even if the introduction discusses histopathology fairness broadly.
  
  \item \textbf{Example E2}: A paper reporting breast cancer classification accuracy across different scanners but without explicit fairness metrics, subgroup analysis, or bias discussion would be excluded under E2.
  
  \item \textbf{Example F2}: A study that uses the term ``bias'' only when discussing bias terms in neural network layers, with no discussion of dataset bias, demographic bias, or domain shift, would be excluded under F2.
  
  \item \textbf{Example F3}: If two papers both report results on TCGA-BRCA using the same model architecture and train/test splits, with Paper B being an extended version of Paper A, only Paper B would be included.
\end{itemize}

% ----------------------------------------------------------------
% Risk of Bias Assessment
% ----------------------------------------------------------------
\section{Risk of Bias Assessment: Adapted QUADAS-2 Framework}
\label{app:methods:quadas2}

We adapted the QUADAS-2 tool for quality assessment of diagnostic accuracy studies, informed by TRIPOD+AI and PROBAST-AI guidelines. Each study was assessed across four domains with signalling questions tailored to AI-based histopathology fairness studies.

\subsection*{Domain 1: Patient Selection}

\textbf{Signalling questions:}
\begin{enumerate}
  \item Was a consecutive or random sample of patients/specimens enrolled?
  \item Was a case--control design avoided?
  \item Did the study avoid inappropriate exclusions (e.g., excluding low-quality slides without pre-specified criteria)?
  \item Was the cohort representative of the target population (e.g., multi-site, diverse demographics)?
\end{enumerate}

\textbf{Risk of bias judgement:}
\begin{itemize}
  \item \textbf{Low risk}: Consecutive/random sampling, no case--control design, minimal inappropriate exclusions, multi-site or demographically diverse cohort.
  \item \textbf{Unclear risk}: Insufficient detail on sampling method or exclusions.
  \item \textbf{High risk}: Case--control design, convenience sampling, single-site cohort without external validation, or inappropriate exclusions.
\end{itemize}

\subsection*{Domain 2: Index Test (AI Model)}

\textbf{Signalling questions:}
\begin{enumerate}
  \item Were the index test results interpreted without knowledge of the reference standard (i.e., was the test set truly held-out)?
  \item If a threshold was used, was it pre-specified rather than optimized on the test set?
  \item Was the model development process described in sufficient detail for replication?
  \item Was the model trained and evaluated using appropriate cross-validation (e.g., patient-level rather than slide-level splitting)?
\end{enumerate}

\textbf{Risk of bias judgement:}
\begin{itemize}
  \item \textbf{Low risk}: Held-out test set not used for training or threshold selection, pre-specified thresholds, detailed methodology, patient-level cross-validation.
  \item \textbf{Unclear risk}: Insufficient detail on train/test separation or threshold selection.
  \item \textbf{High risk}: Data leakage (e.g., slide-level CV when multiple slides per patient), test set used for model selection, or insufficient methodological detail.
\end{itemize}

\subsection*{Domain 3: Reference Standard}

\textbf{Signalling questions:}
\begin{enumerate}
  \item Is the reference standard likely to correctly classify the target condition (e.g., expert pathologist review, validated clinical assay)?
  \item Was the reference standard interpreted without knowledge of the index test results?
  \item For biomarker or molecular prediction tasks, was the reference standard a clinically validated assay (e.g., IHC, NGS)?
\end{enumerate}

\textbf{Risk of bias judgement:}
\begin{itemize}
  \item \textbf{Low risk}: Expert pathologist consensus or validated clinical assay as reference standard, blinded interpretation.
  \item \textbf{Unclear risk}: Insufficient detail on reference standard or blinding.
  \item \textbf{High risk}: Non-expert annotation, unvalidated reference standard, or reference standard potentially influenced by index test results.
\end{itemize}

\subsection*{Domain 4: Flow and Timing}

\textbf{Signalling questions:}
\begin{enumerate}
  \item Was there an appropriate interval between index test and reference standard (i.e., same slides/samples)?
  \item Did all patients/slides receive the same reference standard?
  \item Were all patients/slides included in the analysis (i.e., no post-hoc exclusions)?
  \item For multi-site studies, were all sites included in both training and evaluation?
\end{enumerate}

\textbf{Risk of bias judgement:}
\begin{itemize}
  \item \textbf{Low risk}: Same samples for index test and reference standard, consistent reference standard across all samples, all enrolled samples included in analysis, all sites represented in evaluation.
  \item \textbf{Unclear risk}: Insufficient detail on sample flow or exclusions.
  \item \textbf{High risk}: Different samples for index test and reference standard, inconsistent reference standards, substantial post-hoc exclusions, or sites excluded from evaluation.
\end{itemize}

\subsection*{Applicability Concerns}

In addition to risk of bias, we assessed applicability concerns for each domain:

\begin{itemize}
  \item \textbf{Patient Selection}: Does the study cohort match the review question (histopathology AI fairness)?
  \item \textbf{Index Test}: Does the model architecture and task match the review question?
  \item \textbf{Reference Standard}: Is the reference standard appropriate for the claimed clinical application?
\end{itemize}

\textbf{Applicability concern judgement:}
\begin{itemize}
  \item \textbf{Low concern}: Study directly addresses the review question with appropriate cohort, methods, and reference standard.
  \item \textbf{Unclear concern}: Insufficient detail to assess applicability.
  \item \textbf{High concern}: Study addresses a different question than the review (e.g., technical bias without fairness implications) or uses inappropriate methods for the claimed application.
\end{itemize}

% ----------------------------------------------------------------
% Inter-Rater Reliability
% ----------------------------------------------------------------
\section{Inter-Rater Reliability}
\label{app:methods:reliability}

Title and abstract screening was performed by two independent reviewers with 10\% of records double-screened to assess inter-rater reliability. 

\begin{itemize}
  \item \textbf{Number of double-screened records}: 285 (10\% of 2,847 unique records)
  \item \textbf{Agreement on inclusion/exclusion}: 268/285 (94.0\%)
  \item \textbf{Cohen's $\kappa$}: 0.82 (95\% CI: 0.75--0.89), indicating substantial agreement
  \item \textbf{Disagreements resolved}: 17 records were discussed and resolved by consensus
\end{itemize}

Data extraction was performed by one reviewer with a second reviewer auditing a 20\% random sample ($n=18$ studies):

\begin{itemize}
  \item \textbf{Agreement on bias category assignment}: 16/18 (88.9\%)
  \item \textbf{Agreement on mitigation method classification}: 17/18 (94.4\%)
  \item \textbf{Agreement on metric reporting}: 18/18 (100\%)
  \item \textbf{Cohen's $\kappa$ for bias categories}: 0.85 (95\% CI: 0.71--0.99)
\end{itemize}

Disagreements were resolved through discussion, and the extraction schema was refined to clarify ambiguous category boundaries.

% ----------------------------------------------------------------
% Data Extraction Schema
% ----------------------------------------------------------------
\section{Data Extraction Schema}
\label{app:methods:schema}

Data were extracted into a structured JSON schema with the following fields:

\begin{verbatim}
{
  "study_id": "unique identifier",
  "bibliographic": {
    "title": "paper title",
    "authors": ["author1", "author2"],
    "year": 2024,
    "journal": "journal name",
    "doi": "10.xxxx/xxxxx"
  },
  "cohort": {
    "cancer_type": "breast/prostate/etc",
    "datasets_used": ["TCGA-BRCA", "CAMELYON16"],
    "num_sites": 3,
    "num_patients": 500,
    "num_slides": 1200,
    "demographic_fields_available": ["age", "sex", "race"]
  },
  "bias_analysis": {
    "bias_categories": ["demographic", "institutional"],
    "fairness_metrics_reported": ["subgroup_auc", "worst_group_acc"],
    "subgroups_analyzed": ["race", "sex", "site"]
  },
  "mitigation": {
    "method_family": "harmonization",
    "specific_method": "ComBat",
    "stage": "feature",
    "effectiveness": "strong"
  },
  "results": {
    "overall_auc": 0.85,
    "subgroup_auc_gap": 0.08,
    "worst_group_auc": 0.77
  },
  "reproducibility": {
    "code_available": true,
    "code_url": "https://github.com/...",
    "model_weights_available": false
  }
}
\end{verbatim}

The complete JSON extraction schema is available as Supplementary File~\ref{app:schema_json}.

% ----------------------------------------------------------------
% Excluded Non-English Studies
% ----------------------------------------------------------------
\section{Excluded Non-English Studies}
\label{app:methods:excluded_languages}

Our search was limited to English-language publications. During screening, we identified 23 potentially relevant non-English studies (17 Chinese, 4 Spanish, 2 German). Abstracts were reviewed when available in English. Common exclusion reasons included:

\begin{itemize}
  \item Focus on technical image processing without fairness analysis
  \item Radiology or cytology rather than histopathology
  \item Review articles without original data
\end{itemize}

A complete list of excluded non-English studies with translated titles is provided in Supplementary File~\ref{app:excluded_languages}.
